\section{Data Treatment and Normalization}
\label{sec:data_treatment}

As discussed in the previous section, most of the data registered in the reports concerns a player's jersey number and name.
This poses two problems: first, some clubs adopt a variable jersey number distribution during the season, where players in the starting line-up always have numbers from 1 to 11.
Secondly, using player names can be problematic because we have both the player's nickname and ``full name'', but the full name is often incomplete or truncated.
To address these issues, we convert all player references to CBF's unique ID system.
During the data extraction process, we substitute the player's name and jersey number with their personal CBF ID, ensuring consistent identification across matches and seasons.

Furthermore, our data processing workflow segments each game using all player substitutions and red cards to create temporal segments.
For each game, we generate ``sub-games'', each representing a period during which the same set of players was on the pitch, bounded by substitutions or red card events.
The final result is a JSON file with all games structured as shown in Figure~\ref{fig:json_structure}.
\begin{center}
    \lstset{style=json}
    \begin{lstlisting}
{
    "001": {
        "Summary": {
            "Home": "Vasco da Gama / RJ",
            "Away": "Portuguesa / SP",
            "Result": "1 X 0",
            "Date": "25/05/2013",
            "Time": "18:30",
            "Stadium": "São Januário / Rio de Janeiro"
        },
        "0": {
            "Home": {
                "Squad": ["146379", "186825", "430895", "165075", "402413", "307406",
                          "357424", "317751", "156194", "167111", "175011"],
                "Cards": [[[28, False], "430895"]],
                "Goals": [[[47, False], "402413"]]
            },
            "Away": {
                "Squad": ["159483", "158204", "160789", "134212", "151104", "160357",
                          "176908", "138048", "190258", "136911", "188158"],
                "Cards": [[[23, False], "176908"]],
                "Goals": []
            },
            "Time": 62
        },
        "1": {
            "Home": {
                "Squad": ["146379", "186825", "430895", "165075", "402413", "307406",
                          "357424", "317751", "156194", "167111", "143841"],
                "Cards": [],
                "Goals": []
            },
            "Away": {
                "Squad": ["159483", "158204", "160789", "134212", "151104", "160357",
                          "176908", "138048", "190258", "136911", "188158"],
                "Cards": [],
                "Goals": []
            },
            "Time": 4
        }
    }
}
    \end{lstlisting}
    \captionof{figure}{
        \textbf{Final JSON structure example}.
        Example of the game between CR Vasco da Gama and Associação Portuguesa de Desportos, showing two sub-games triggered by a substitution (CBF ID 175011 replaced by 143841) at the 62nd minute.
        The ``Summary'' field stores general metadata such as team names, final result, date/time, and venue.
    }
    \label{fig:json_structure}
\end{center}

The JSON structure can be interpreted as follows:
\begin{itemize}
    \item \texttt{"001"}: Represents the unique identifier for the match.
    \begin{itemize}
        \item \texttt{"Summary"}: Provides a summary of the match, including home team, away team, final result, date, time, and stadium.
        \item Other second-level keys (e.g., \texttt{"0"}, \texttt{"1"}, \texttt{"2"}, ...) represent different segments of the match.
        These segments are created whenever a substitution or red card occurs, even if the red card is shown to a reserve player.
        \begin{itemize}
            \item \texttt{"Home"} and \texttt{"Away"}: Information about each team during this segment.
            \begin{itemize}
                \item \texttt{"Squad"}: List of player CBF IDs for the team.
                \item \texttt{"Cards"}: List of cards received by team players.
                Each card is represented by an array where the first element is a tuple with the minute the card was shown and a boolean flag indicating whether it occurred during injury time of the current segment.
                The second element is the player's CBF ID.
                \item \texttt{"Goals"}: List of goals scored by the team, each represented by an array where the first element is a tuple with the minute the goal was scored and a boolean flag indicating whether it occurred during injury time.
                The second element is the scorer's CBF ID.
            \end{itemize}
            \item \texttt{"Time"}: Duration of this segment in minutes.
        \end{itemize}
    \end{itemize}
\end{itemize}

An important consideration in this process is that all reports are written manually, which means they may contain typos or mistakes.
When we identified anomalies that caused errors in the retrieval process, we manually corrected them.
However, this correction process was only performed for First Division (\textit{Série A}) matches.
Consequently, in lower divisions, we can expect more anomalies that were not treated, which may result in potentially inferior data quality for those competitions.
